\documentclass[11pt,a4paper]{article}
\usepackage[T1]{fontenc}
\usepackage[utf8]{inputenc}
\usepackage[czech]{babel}
\usepackage{lmodern,microtype}
\usepackage{amsmath,amssymb,mathtools}
\usepackage{geometry}
\geometry{left=19mm,right=19mm,top=21mm,bottom=21mm}
\usepackage[dvipsnames]{xcolor}
\usepackage{enumitem,tabularx,booktabs}
\usepackage[most]{tcolorbox}
\usepackage{fancyhdr,lastpage}
\usepackage[hidelinks]{hyperref}
\definecolor{deep}{HTML}{203955}
\definecolor{teal}{HTML}{176E76}
\definecolor{slate}{HTML}{506173}
\definecolor{pale}{HTML}{F3F7F9}
\setlength{\parindent}{0pt}
\setlength{\parskip}{5pt}
\setlist[enumerate]{leftmargin=6mm,itemsep=5pt,topsep=3pt}
\pagestyle{fancy}\fancyhf{}
\fancyhead[L]{\sffamily\footnotesize\color{slate} MAT · 4. ročník · kombinatorika}
\fancyhead[R]{\sffamily\footnotesize\color{slate} Hodina \LESSONNUM\, / 112}
\fancyfoot[L]{\sffamily\footnotesize\color{slate} Výukový podklad · samostatný dokument}
\fancyfoot[R]{\sffamily\footnotesize\color{slate} \thepage/\pageref{LastPage}}
\renewcommand{\headrulewidth}{0.35pt}
\renewcommand{\footrulewidth}{0pt}
\newcommand{\sect}[1]{\par\vspace{2mm}{\Large\sffamily\bfseries\color{deep} #1}\par\vspace{-0.5mm}\noindent\color{teal}\rule{\linewidth}{0.7pt}\color{black}\vspace{1mm}}
\tcbset{theobox/.style={enhanced,breakable,colback=pale,colframe=teal!45!white,boxrule=.6pt,arc=1mm,left=3mm,right=3mm,top=2mm,bottom=2mm}}
\newcommand{\LESSONNUM}{40}
\begin{document}

{\Huge\sffamily\bfseries\color{deep} Průběžný test: variace, permutace a kombinace}\par
{\sffamily\large\color{teal} Výukový podklad · 40. vyučovací hodina}\par
\textbf{Hodinová dotace:} 45 minut\quad \textbf{Ročník:} 4. SPŠ\quad \textbf{Rozpis:} TP T3\par
\begin{tcolorbox}[theobox]\textbf{Výukový cíl.} Získat diagnostický podklad pro 15–20minutové ověření základních kombinatorických postupů a pro následnou opravu typických chyb.\end{tcolorbox}
\sect{Teorie a vysvětlení}
\textbf{Tato hodina je obsahový a metodický podklad, nikoli schválený žákovský test.} Vlastní test musí být později vytvořen pomocí projektu \texttt{vytvor-test} a výběr úloh musí schválit učitel.\par
\textbf{Rozsah:} pravidlo součinu, variace bez opakování a s opakováním, permutace bez opakování a s opakováním, kombinace bez opakování. Vzhledem k názvu testu v rozpisu nezavádíme nové učivo a binomickou větu nepovažujeme za povinnou testovou položku.\par
\textbf{Organizace 45 min:} \(0\)--\(20\) min krátký test; \(20\)--\(27\) min porovnání strategií bez zveřejňování individuálních výsledků; \(27\)--\(40\) min opravné úlohy; \(40\)--\(45\) min shrnutí rozhodovacího algoritmu.
\sect{Řešené příklady}
\textbf{Příklad 1.} Diagnostika A: Kolik různých trojciferných čísel lze napsat z \(0,1,2,3,4,5\), bez opakování?\par
\textit{Řešení.} První cifra má \(5\) nenulových možností; druhá \(5\) a třetí \(4\). Počet \(5\cdot5\cdot4=100\). Běžná chyba je \(6\cdot5\cdot4=120\), která povoluje úvodní nulu.\par\smallskip
\textbf{Příklad 2.} Diagnostika B: Kolik různých pořadí písmen vznikne ze slova \textsc{AABBC}?\par
\textit{Řešení.} Jde o permutace s opakováním: \(\frac{5!}{2!\,2!}=30\). Chybný výsledek \(120\) přehlíží nerozlišitelné kopie.\par\smallskip
\textbf{Příklad 3.} Diagnostika C: Z osmi osob vyber tři členy nerozlišené komise.\par
\textit{Řešení.} Na pořadí nezáleží. \(\binom83=56\). Číslo \(8\cdot7\cdot6=336\) je počet uspořádaných trojic, šestkrát větší.\par\smallskip
\textbf{Příklad 4.} Diagnostika D: Ze šesti družstev urči vítěze, druhého a třetího.\par
\textit{Řešení.} Jde o variace bez opakování: \(6\cdot5\cdot4=120\).\par\smallskip
\textbf{Příklad 5.} Diagnostika E: Kolik čtyřznakových PINů z číslic \(0\) až \(9\) je možných?\par
\textit{Řešení.} Variace s opakováním: \(10^4=10\,000\). PIN nemá zákaz počáteční nuly.\par\smallskip
\clearpage\sect{Procvičování pro studenty}
\begin{enumerate}
\item Oprava 1: Tříznakový kód ze čtyř symbolů s opakováním.
\item Oprava 2: Všechny možné řady pěti různých osob.
\item Oprava 3: Z deseti lidí vyber dva bez rozlišení funkcí.
\item Oprava 4: Uspořádej písmena \textsc{MAMA}.
\item Oprava 5: Čtyřciferná čísla z cifer \(0\) až \(4\), bez opakování.
\end{enumerate}
\sect{Řešení a zdůvodnění k procvičování}
\begin{enumerate}
\item \(4^3=64\).
\item \(5!=120\).
\item \(\binom{10}{2}=45\).
\item \(\frac{4!}{2!2!}=6\).
\item \(4\) možnosti první cifry, pak \(4,3,2\): \(4\cdot4\cdot3\cdot2=96\).
\end{enumerate}
\sect{Metodické poznámky}
\textbf{Typické chyby:}
\begin{itemize}[leftmargin=5mm,itemsep=1pt]
\item Nikdy nevyvozovat z jediného výsledku, že student pochopil typ úlohy; žádat stručné slovní zdůvodnění.
\item Při opravě zdůraznit nepřípustnou počáteční nulu, dvojnásobné započítání a nerozlišitelné prvky.
\item Jde o obsahový návrh, nikoli finální písemnou práci pro studenty.
\end{itemize}
\textbf{Doporučený průběh 45 minut:} 0–20 diagnostický test; 20–27 rozbor strategií; 27–40 cílené opravy; 40–45 pravidla rozhodování.\par
'\textbf{Doporučená diagnostika:} Vybrat čtyři až pět základních položek, celkem přibližně 10--12 bodů. Každá musí ověřovat odlišný typ počítání. Nepředkládat jako hotový test bez učitelova schválení.\par
\textbf{Návaznost:} Po diagnostice následují smíšené úlohy s omezeními a most k pravděpodobnosti.\par
\vspace{2mm}\noindent{\footnotesize\color{slate} \textbf{Východiska:} rozpis výuky MAT, 4. ročník, hodiny 29--41; tematický plán SPŠ; srovnání tematického rozsahu a náročnosti s realisticky.cz, oddíl Kombinatorika (lekce 090101--090119). Výklad, příklady a řešení jsou nově autorsky sestavené, nikoli převzatá sbírka.}\par
\end{document}